\section{CausalPy at a lending threshold: what a reusable experiment must preserve}
\label{exec14:causalpy-contract}

A credit policy often turns a continuous score into a discrete action. A borrower
just above a threshold might receive a different offer, servicing contact or
review requirement. The useful software question is therefore not simply whether
a package can fit two curves. It is whether the resulting contrast answers the
bank's intervention question, whether its uncertainty refers to that contrast,
and whether the interface preserves the design when observations, units or models
change. CausalPy offers an instructive candidate because its experiment interface
can expose a Bayesian regression or a point-estimation backend. The source
revision considered here declares version 1.0.0rc1 and is dated 6 October 2026;
the earlier distribution-metadata observation of 0.9.0 is a different artifact.
The following analysis concerns the selected sharp regression-discontinuity
path, not automatic approval of its other designs.\citep{exec14_causalpy_primary}

\subsection{The local question precedes the fitted contrast}

Let $X$ be the score available at the decision time, $c$ the cutoff and
$T=\mathbf{1}\{X\geq c\}$ the stipulated sharp assignment rule. Define $Y(t)$ as
the fixed-horizon outcome under action $t$. If observed outcomes obey consistency,
the conditional potential-outcome means are continuous at $c$, and score support
approaches the cutoff from both sides, the limiting observed jump identifies
the local mean effect. To see the dependence on these assumptions explicitly,
write $\mu_t(x)=E[Y(t)\mid X=x]$. Sharp assignment and consistency imply
$E[Y\mid X=x]=\mu_0(x)$ below the cutoff and $\mu_1(x)$ above it. Continuity
then gives
\[
 \tau_c=\mu_1(c)-\mu_0(c)
 =\lim_{x\downarrow c}E[Y\mid X=x]-\lim_{x\uparrow c}E[Y\mid X=x].
\]
This derivation is conditional on the stipulated design. It does not establish
that a real score cannot be manipulated, that a simultaneous pricing change is
absent, or that the effect transports to all borrowers. A servicing outcome also
requires a stable measurement horizon and a declared treatment of censoring.
Those are parts of the question, not properties acquired by importing a regression
class.\citep{exec14_causalpy_exact,exec14_causalpy_decision}

The inspected CausalPy experiment builds design matrices and later obtains
predictions at two rows: $(c-\epsilon,T=0)$ and $(c+\epsilon,T=1)$. Its result
subtracts the first prediction from the second. The default separation is
$\epsilon=0.001$ in the running variable's units. Consequently the returned
quantity is a finite-separation model contrast,
\[
 \Delta_\epsilon=m(c+\epsilon,1)-m(c-\epsilon,0),
\]
which approaches a model's threshold jump only under the relevant limiting
behavior. The implementation's name does not make a finite separation identical
to that limit. The inspected treatment rule assumes treatment above the cutoff;
the selected validation method requires a dummy-coded treated column but does
not itself demonstrate that every observed treatment matches this rule.
An eligibility cutoff with imperfect take-up is a different design. Relabeling
actual receipt as deterministic eligibility would hide the first-stage question.
\citep{exec14_causalpy_primary}

Consider the exact, stipulated conditional mean with the score centered at zero,
\[
 m(x,t)=100+20x+5t+10xt.
\]
Its threshold jump is 5. At $\epsilon=1/1000$, however,
$\Delta_\epsilon=5+(2\cdot20+10)/1000=5.05$.
More generally, for $m(x,t)=\alpha+\beta x+\tau t+\gamma xt$ centered at the cutoff,
the difference is $(2\beta+\gamma)\epsilon$. This is an exact contrast identity,
not an observed software bias estimate or an assertion about a fitted sample.
Shrinking the separation reduces this particular deterministic difference but
does not repair a misspecified regression, weak support or an invalid design.
\citep{exec14_causalpy_exact}

Units make the distinction operational. Expressing the same score as $Z=100X$
changes the coefficients to $0.2$ and $0.1$ for the score and its interaction.
Keeping the numeric separation at 0.001 now produces 5.0005; preserving the
physical separation requires $\epsilon_Z=0.1$, which restores 5.05. The cutoff,
bandwidth and donut exclusion must also transform. A unit conversion that leaves
all numeric tuning inputs untouched therefore changes the question asked of an
otherwise equivalent model. A proposed adapter should return the two physical
prediction locations and their units alongside the contrast.
\citep{exec14_causalpy_exact,exec14_causalpy_decision}

\subsection{What fit quality cannot resolve}

The need for a design restriction can be demonstrated without simulated noise.
In world A let $Y(t)=100+20X+5t$; in world B let
$Y(t)=100+20X+5\mathbf{1}\{X\geq0\}$ for both treatment values. Retain
$T=\mathbf{1}\{X\geq0\}$ in each world. Both worlds have exactly the same
observed outcome function, including the jump of 5, but their treatment effects
are respectively 5 and 0. World B violates the continuity restriction. No
estimator that sees only this observed law can use fit quality to decide which
counterfactual extension is true. For an illustrative outcome measured in dollars
and an action cost of USD 3 per borrower near the cutoff, the stipulated net
effect is USD 2 in world A and minus USD 3 in world B. These are alternative
local mechanisms, not a bank-wide return estimate.\citep{exec14_causalpy_exact}

This ambiguity suggests a useful public-data prototype. Generate known mechanisms
with and without baseline discontinuities, retain the same observed law where
possible, and require the workflow to state which assumption separates them.
Public lending or policy datasets can exercise data preparation and diagnostics,
but their public availability alone does not reveal the missing counterfactual
truth. Density or covariate checks can contribute design evidence; neither should
be represented as a proof of every identifying restriction. The evaluation should
score appropriate qualification and refusal as well as numerical recovery.
The proposed prototype can therefore demonstrate whether the interface preserves
the question before an expensive institution-specific study is justified.
\citep{exec14_causalpy_decision}

\subsection{Uncertainty is a semantic type, not an array dimension}

The inspected Bayesian linear model places Normal priors with scale 50 on
coefficients and uses a Normal likelihood whose noise scale has a HalfNormal
prior with scale 1. Its deterministic mean is the design matrix times the
coefficient vector. The experiment adapter extracts these expected-outcome means;
the response-predictive variable has a different interpretation. The selected RD
plotting path uses equal-tail posterior quantiles. Such quantiles condition on the
model and prior: they are not automatically highest-density intervals, robust
bias-corrected frequentist intervals, or uncertainty about the truth of the causal
design. The sampling path merges stored settings and call-time overrides before
invoking the sampler. A reproducible request consequently needs the actual runtime
and sampler configuration, not just the experiment class name.
\citep{exec14_causalpy_primary}

A point-estimation backend introduces a subtler conversion hazard. Its adapter
places point predictions into arrays with singleton chain and draw dimensions.
Those dimensions provide a common representation; they do not manufacture
posterior variation. An interval routine applied mechanically to a singleton
would return equal endpoints, but that degeneracy is not evidence of statistical
certainty. The proposed result contract must distinguish a point estimate from
posterior expected outcomes, response-predictive draws and prior draws. Requests
for unsupported uncertainty should return an explicit unavailable result, never
a zero-width confidence claim.\citep{exec14_causalpy_primary,exec14_causalpy_decision}

Prior scales must follow units as well. In the coordinate change $Z=100X$, a
coefficient distributed as Normal with scale 50 in the old score units has scale
0.5 in the new units if the prior is to describe the same model. Resetting both
to the same default numeric scale changes prior beliefs. Outcome transformations
likewise affect coefficient and noise priors. These observations do not condemn
the default model: they explain why a reusable banking operation must record and
review its units, likelihood and prior rather than silently rebuilding them.
\citep{exec14_causalpy_primary,exec14_causalpy_exact}

\subsection{Support, lifecycle and the choice of an alternative}

The selected implementation filters a symmetric bandwidth and optionally removes
a donut around the threshold. At ten or fewer retained observations it warns;
the warning is not a scientifically sufficient qualification rule. A positive
donut leaves no observations arbitrarily close to the cutoff in that selected
sample, so threshold predictions rely on the model's extrapolation across the
excluded region. The default infinite bandwidth also does not constitute a
data-driven local bandwidth choice. A bank wrapper needs left and right support,
row-selection records and an explicit extrapolation flag alongside fitted output.
It should separately check invalid tuning inputs and reversed treatment coding.
These are proposed boundary obligations; the selected source has not been run
against those cases here.\citep{exec14_causalpy_primary,exec14_causalpy_decision}

Lifecycle is equally consequential for an agent or a long-running service. The
source warns that its frozen design fingerprint does not detect later edits to
the experiment's data and recommends constructing a new experiment for changed
data. Assigning a fresh model resets result state. Therefore a request should bind
an immutable data identity and a formula/design identity, create a fresh experiment
when either changes, and keep configured, built and fitted states distinguishable.
The selected test sources cover RD smoke paths, treatment normalization and donut
masks. Reading those tests is evidence of intended cases, not evidence that this
proposal executed them or obtained calibrated RD intervals.
\citep{exec14_causalpy_primary}

For a different uncertainty objective, the inspected R documentation for
\texttt{rdrobust} describes local-polynomial estimators and robust bias-corrected
confidence intervals, with separate fitting and bias bandwidths, polynomial
orders, kernels and sharp/fuzzy/kink options. This supplies a concrete comparison
to a Bayesian model contrast. The choice depends on the design, target and
uncertainty statement; the packages should not be interchanged by a generic
causal-effect endpoint. The R source snapshot used for this comparison
precedes the stated literature cutoff. Documentation of an algorithm is not a
qualification result for the bank's sampling law.\citep{v12_sw_rdrobust}

\subsection{A scoped reuse decision and the work it creates}

The proposal retains CausalPy as a pinned reference and a candidate for a narrow
future integration. Direct production adoption is deferred until the conversion
and failure contract has been implemented and exercised. Porting the model or
sampler is not presently justified: it transfers numerical and maintenance
ownership without an established benefit. Rejecting the entire package is also
unsupported by this review. Its root source license is Apache 2.0, its packaging
requires Python 3.12 or later, and its declared direct dependencies include broad
PyMC and PyTensor major-version ranges. These declarations establish neither a
resolved environment nor complete dependency rights. The pinned pre-release
version and dated activity are observable facts; they are not promises of future
support.\citep{exec14_causalpy_primary,exec14_causalpy_decision}

The next prototype should exercise cutoff equality, direction reversals, physical
unit transformations, bandwidth and donut masks, changed-data reconstruction,
point/posterior/predictive distinctions and failure retention. A scientific owner
must specify the local estimand, continuity restrictions and the uncertainty
claim. An integration owner must preserve data, formula and prediction identities.
A runtime owner must resolve dependencies, retain notices, fix sampler settings
and define diagnostic failures. An evaluation owner must predeclare generated
mechanisms, public-data uses and appropriate recovery and refusal criteria.
DiD, synthetic control, IV and longitudinal uses remain separate design decisions.
\citep{exec14_causalpy_decision}

Costs should include environment resolution, model and prior review, sampling
time, failed fits, diagnostic inspection, conversion maintenance and comparison
against a simpler alternative. These costs can be measured on public or generated
cases before committing to bank experiments. The evidence supports a provisional product choice and specifies the tests needed before integration. Integration
remains disabled until executed upstream behavior supports the proposed contract;
the exact examples in this section establish mathematical distinctions, not a
working external adapter.\citep{exec14_causalpy_decision}
\label{exec14:causalpy-contract-end}
